\DocumentMetadata{
  pdfversion=2.0,
  lang=es-MX,
  pdfstandard=ua-2
}

\documentclass{sener2025}

% --- Metadatos del Documento ---
\title{Prueba Normalización de Saltos}
\subtitle{Verificación de corrección "There's no line here to end"}
\author{Equipo de Desarrollo}
\date{13 de diciembre de 2025}
\institucion{Secretaría de Energía}
\unidad{Dirección de Sistemas}
\setDocumentoCorto{PRUEBA\_SALTOS}
\palabrasclave{SENER, LaTeX, Saltos, Normalización}
\version{1.0}

\begin{document}

\portadafondo

\tableofcontents
\newpage

\section{Pruebas de Normalización de Saltos}

\subsection{Caso 1: Texto normal sin problemas}

Este es un texto normal que no debería causar problemas. Tiene párrafos separados correctamente.

Este es otro párrafo que viene después de una línea en blanco.

\subsection{Caso 2: Resumen Ejecutivo (caso problemático original)}

\begin{resumenejecutivo}
Objetivo: detectar errores de compilación causados por saltos de línea mal procesados.

Este texto simula el contenido que venía de Google Sheets con saltos problemáticos.

Antes se generaban comandos \textbackslash\textbackslash{} al inicio de párrafos causando "There's no line here to end".
\end{resumenejecutivo}

\subsection{Caso 3: Recuadro con saltos internos}

\begin{recuadro}{Prueba de Saltos}
Primera línea del recuadro.
Segunda línea que debería estar en el mismo párrafo.

Esta línea debería iniciar un nuevo párrafo dentro del recuadro.
\end{recuadro}

\subsection{Caso 4: Lista con elementos multilínea}

\begin{itemize}
  \item Primer elemento que puede tener texto largo y saltos internos sin causar problemas
  \item Segundo elemento con contenido normal
  \item Tercer elemento que también funciona correctamente
\end{itemize}

\section{Verificación de Corrección}

Si la normalización funciona correctamente:
\begin{itemize}
  \item ✅ No aparece el error "There's no line here to end"
  \item ✅ Los párrafos se separan correctamente
  \item ✅ Los saltos simples se convierten en espacios
  \item ✅ Los saltos dobles se convierten en \textbackslash par
  \item ✅ No hay \textbackslash\textbackslash{} al inicio de bloques
\end{itemize}

\end{document}